\section{Agent Swarm: colaboración multiagente}

\subsection{El orquestador y la organización de subagents}
El núcleo de esta parte del entrenamiento no son los subagents en sí, sino el aprendizaje del Allcastrater (el orquestador y descompositor de tareas); el instructor enfatiza «lo que entrenamos es al descompositor orquestador». El Allcastrater se encarga de invocar \texttt{createSubagents} y \texttt{assignTask}, asignar tareas conforme a distintos system prompts, recolectar resultados y luego seguir generando nuevos subagents. Su conducta sincroniza información entre cada ronda, orquesta subagents en paralelo dentro de los ciclos internos, y observa continuamente la reward que entrega el GRM para ajustar la scheduling policy.

\begin{knowledgebox}{Allcastrater: un orquestador que se puede aprender}
\begin{itemize}
    \item El Allcastrater construye múltiples system prompts según la tarea y los recursos actuales, garantizando que cada subagent se concentre en una subtarea específica (por ejemplo, razonamiento sobre imágenes, generación de texto, invocación de herramientas).
    \item Usa \texttt{createSubagents} para generar nuevos agents, \texttt{assignTask} para asignar el trabajo y \texttt{collectResults} para combinar las respuestas, manteniendo la continuidad de la tarea previa.
    \item El objetivo de entrenamiento del Allcastrater es la calidad de la orquestación y no el detalle de cada subagent, de modo que el sistema completo se puede evaluar de manera unificada mediante el GRM y ajustar dinámicamente el compute budget dentro del parallel loop.
\end{itemize}
\end{knowledgebox}

\subsection{Descomposición de tareas y flujo de herramientas}
Durante el entrenamiento, el Allcastrater fragmenta las tareas cada vez más: primero identifica el high-level intent (por ejemplo, «reconstruir la página principal del Apple iPad Air»), y después usa distintos subagents para generar el análisis, el diseño, los bloques de código y las pruebas. Cada subagent invoca la herramienta correspondiente (renderizador de imágenes, editor de HTML, consulta a base de datos) y devuelve el resultado al scheduler, que decide la siguiente revisión del prompt. Como cada subagent tiene su propia prompt template, también hay que registrar los parameters y el timestamp de cada tool calling.

\begin{knowledgebox}{Las tres dimensiones del tool pipeline}
\begin{itemize}
    \item Agent de análisis: lee los requisitos, valida las restricciones del negocio, decide el tipo de subagent que sigue, y transmite el intent al scheduler como JSON estructurado.
    \item Agent de generación: construye texto explicativo, llamadas a API, fragmentos de código e incluso mockups de interfaz.
    \item Agent de verificación: usa GUI agents o CLI agents para leer la retroalimentación de las herramientas, determina si se cumplen los critical steps y reporta el estado al Allcastrater.
\end{itemize}
\end{knowledgebox}

\subsection{Caso de evaluación de agent swarm}
En la evaluation list que presenta el instructor aparecen tareas como replicating Apple iPad Air site, descripción de video, operaciones de GUI, entre otras; todas se dividen en varios subagents: un agent visual que extrae la información de la interfaz, un agent de texto que redacta las explicaciones y un agent de herramientas que ejecuta las API. Al terminar cada agent loop, el Allcastrater usa el GRM para calificar el conjunto completo de respuestas y decide si conviene seguir generando (spawn) nuevos subagents para pulir aún más la salida; si el GRM detecta que cierto tipo de tool calling tiene una reward persistentemente baja, esa action se suspende en el siguiente loop hasta ajustar la prompt grammar.

\subsection{Rondas en paralelo y S mean + max}
Dentro de cada ronda, los subagents se ejecutan en paralelo; entre rondas, en cambio, la sincronización sigue estrictamente al Allcastrater entrenable: al completar un lote de tareas se recolectan los resultados, y después se asignan nuevas tareas o se genera (spawn) un nuevo subagent. Como muestra el instructor en el monitoreo del entrenamiento, S mean + max se usa para agregar el desempeño de una ronda, y el conjunto completo del agent swarm exhibe una eficiencia de critical steps 4.5 veces mayor que la de un solo agent.

\begin{importantbox}{Eficiencia de 4.5 veces y la Vision-Agentic Capability}
\textit{``Activate the Vision-Agentic Capability''} no es una consigna, sino la indicación de que cada agent loop debe percibir simultáneamente lo visual, el lenguaje y las llamadas a herramientas antes de decidir el siguiente paso. Esto permite que el conjunto de subagents agregue puntajes mediante \texttt{S mean + max} y complete, con una tasa de ahorro de recursos de 4.5 veces, el número de pasos que originalmente solo un single agent podría alcanzar.
\end{importantbox}

\subsection{Critical Steps y el agent loop}
Para acotar el compute del agent swarm, tanto el entrenamiento como la evaluación usan un umbral de «Critical Steps» que limita el número de veces que se orquesta en cada ronda, manteniendo un estado estable bajo la resource constraint. Al explicar esto, el instructor menciona el VRL agent loop, los GUI agents y el workflow de tool calling: el agent loop primero decide el intent, después invoca las herramientas y por último regresa al Allcastrater. La tarea de los GUI agents es incorporar los pasos de interacción con la interfaz dentro de este loop, de modo que el modelo conserve, en tareas que cruzan distintos dominios, la capacidad de controlar la interfaz y obtener los valores de retorno de las herramientas.

\begin{warningbox}{No trates los critical steps como un presupuesto fijo}
El instructor advierte que los critical steps no son un umbral estático, sino una guard band ajustable: en cuanto el modelo se vuelve estable en cierto benchmark, se puede invertir el budget restante en un nuevo subagent; a la inversa, si el GRM detecta que la reward de toda la ronda baja, se deben reducir los steps de inmediato en lugar de simplemente añadir más tiempo sin criterio.
\end{warningbox}

\subsection{Resumen de la sección}
La clave de agent swarm está en un Allcastrater que se puede aprender, en subagents que corren en parallel y en el mecanismo de control de recursos de los Critical Steps. Basta con mantener cerrado el ciclo de \texttt{createSubagents} \& \texttt{assignTask} y registrar el tool trace a tiempo para que K2.5 pueda producir más agentic output con menos compute.
